\section{Introduction}

International shipping faces simultaneous pressure to reduce greenhouse-gas emissions, diversify its fuel base, and maintain operational reliability.
The IMO 2023 strategy and industry orderbook trends have turned alternative-fuel propulsion from a distant vision into a practical engineering reality
\cite{imo2023ghgstrategy,classnk2025alternativefuels}.
As fuels such as LNG, methanol, and ammonia are introduced to lower carbon intensity, marine dual-fuel and multi-fuel engines have become an important development direction for large propulsion systems
\cite{wingd_lng_faq_2025,wingd_ammonia_faq_2025,classnk_methanol_2024}.

However, alternative fuels introduce further challenges to marine-engine combustion processes.
To enable practical operation with low-reactivity primary fuels, an ignition source with high reactivity and reliability is still required.
The fuel system may become lower-carbon and more flexible, but combustion still depends on a small yet critical local event: pilot diesel injection.
Pilot diesel injection remains a critical enabling mechanism in multi-fuel marine engines because it initiates combustion in a low-reactivity alternative-fuel main charge that does not auto-ignite under practical compression-ignition conditions
\cite{wingd_xdf_manual_2021,classnk_methanol_2024}.
For injector development, one central control problem is the spatial development of the pilot spray itself.
Spray penetration, cone angle, and morphology influence the position and volume of the ignition kernel, and that placement affects main-fuel mixing, flame propagation, and overall combustion efficiency
\cite{Karim2015DualFuel,WeiGeng2016DualFuelReview,Sahoo2009DualFuelReview}.
If pilot penetration is too short, the ignition kernel may remain too small or form in an unfavourable region; if it is too long, the spray may move away from the intended design and increase wall-interaction risk.

When that control objective is not met, the same pilot injection may become a source of operational risk instead.
If its liquid envelope reaches the piston crown or cylinder wall before a useful ignition kernel forms, wall wetting, lube-oil dilution, degraded mixture preparation, and increased unburned-hydrocarbon emissions may follow
\cite{peraza2022swi,ahmad2025fueldilution,MorieiraMotaPanao2010}.
The practical question is whether its spatial development remains inside a controllable window that supports ignition, mixing, and safety margin.

Existing research and engineering practice provide several established approaches for assessing pilot-spray development.
Engine tests examine spray, combustion, wall interaction, and emissions within an integrated system, while high-fidelity CFD can resolve richer flow details.
Empirical penetration correlations and phenomenological models describe spray development rapidly; comparing penetration with wall distance also provides a natural first-order geometric criterion for impingement screening
\cite{HiroyasuArai1990,NaberSiebers1996,Baumgarten2006MixFormation}.
Optical spray experiments occupy a complementary niche by providing repeatable observations of the liquid envelope.
Among them, high-speed Mie imaging is especially suitable for observing liquid-phase boundaries and macroscopic penetration evolution
\cite{Linne2013SprayImaging}.

Each approach, however, covers only part of the need for fast, condition-parametric, experimentally grounded pilot-spray penetration prediction.
Engine tests are expensive and slow; CFD requires complex inputs, model tuning, and substantial computation.
Empirical models are compact and interpretable, but short pilot injections make their use more difficult because the observable window can be dominated by start-of-injection and end-of-injection transients
\cite{ZhouLiYi2021,ZhouLiLaiWei2019}.
Optical experiments provide image-based evidence, but image post-processing still produces discrete observation tables rather than a predictive model for fast design iteration.

This thesis therefore addresses a specific missing intermediate layer.
A multihole nozzle produces plumes that are nominally similar but not fully equivalent.
A fixed field of view turns part of the late-time penetration history into right-censored observations, and raw Mie intensity is affected by illumination, glare, and multiple scattering.
If the visible endpoint of each observed trajectory is treated as the true endpoint, the model can learn ``leaving the image'' as ``physical saturation''.
The missing piece is not a new theory of spray physics, nor an additional isolated diagnostic measurement, but a fast, repeatable, explicitly scoped intermediate predictor step between optical experiment archives and high-cost validation.

Data-driven surrogates offer a feasible direction for this goal.
Recent studies show that neural networks can predict macroscopic spray properties such as penetration and cone angle from operating-condition parameters
\cite{VazKarathanasis2026,liu2020sprayreview}.
Here, multi-campaign Mie recordings acquired under the same optical configuration can be converted into plume-resolved temporal observables and used to train a fast mapping from conditions to penetration trajectories.
The central contribution of this thesis is therefore an auditable workflow that connects imaging data, geometric definitions, censoring treatment, uncertainty representation, and simplified wall-risk screening.

The aim of this thesis is to develop and evaluate an auditable screening workflow that converts a large archive of Mie-scattering spray videos into an early-stage engineering prediction layer between optical experiments and high-cost validation.
To accomplish this aim, the thesis carries out four tasks:
\begin{enumerate}
  \item extract plume-resolved geometric observables from high-speed non-reacting Mie imaging videos;
  \item construct penetration-trajectory targets from finite-field-of-view observations;
  \item train a surrogate that predicts both penetration mean and conditional uncertainty; and
  \item connect the predicted distribution to a simplified representative engine geometry to form a probabilistic wall-impingement screening proxy.
\end{enumerate}

\subsection{Main contributions}

The main contributions are a CUDA-accelerated video-to-observable workflow, a reduced-order fitting strategy for right-censored penetration targets, a multi-stage uncertainty-aware learning framework, and a simplified probabilistic impingement-screening method.
The scope is deliberately limited to non-reacting, non-evaporating pilot diesel spray under controlled pressurised-chamber conditions.
The model is an engineering screening tool within its supported data domain, not a universally reliable predictor for arbitrary unseen conditions or a finished truth model for real engine wall impingement.

\subsection{Thesis structure}

Chapter~2 establishes the literature rationale; Chapter~3 describes the experimental data and image-processing workflow; Chapter~4 describes trajectory targets, surrogate modelling, and probabilistic screening; Chapter~5 reports the empirical results; Chapter~6 discusses what has been validated and where the current scope stops; and Chapter~7 separates validated outcomes, prototype boundaries, and next steps.
